\subsection{调度}
\label{sec:algo_other}

在许多场景中，例如使用因果掩码或可变序列长度（varlen）时，注意力内核天然存在负载不均衡——分配给SM的工作块的主循环长度不同，因为某些工作块比其他块需要更多的加载和MMA操作。此外，我们可以选择SM处理块的顺序，例如通过定义网格坐标的首选线性化方式。抽象掉注意力的任何特定功能，我们可以将关于同构并行处理器的最大完工时间（makespan）最小化的一般结果应用于我们的场景。特别地，在FlashAttention-4中，我们使用经典的最长处理时间优先（LPT）调度~\citep{graham1969}。我们强调，我们应用这一思想的方式适用于所有GPU架构，并且在Hopper GPU上也被验证为FlashAttention-3的改进。

\textbf{因果掩码的LPT调度。}标准注意力网格为（m块，头，批次），按从左到右递增的顺序计算。但是分数在对角线上方被掩码掉，因此对于固定的头和批次，SM最终从最短到最长低效地处理工作块。另一方面，朴素的LPT顺序也非最优，因为对于不同批次，主循环KV加载不会命中L2缓存，而且如果KV头超过L2缓存容量则会发生缓存抖动。相反，我们始终将批次作为最外层维度处理，并对头进行重排。这意味着我们将头分成不超过L2缓存容量的组；分块调度器随后按每组头、逆序的m块、组别、最后批次的顺序遍历网格。特别地，对于MQA或GQA，我们总是在改变m块之前遍历每个KV头的所有查询头。经验上验证了这种LPT顺序非常有效；例如，对于BF16和头维度128，在H200 GPU上测量到MHA获得4--8\%的FLOPS提升，MQA 8获得7--14\%的提升。

\textbf{可变序列长度的LPT调度。}对于varlen，我们还需要处理批次间变化引起的负载不均衡。例如，在解码工作负载中，不同批次可能关注不同数量的上下文；在混合批处理或连续批处理中，某些批次可能处于预填充（prefill）阶段而其他处于解码阶段。每个批次的查询和KV序列长度列表通常作为注意力元数据存储在设备上，标准的varlen注意力内核在以递增顺序处理批次时在运行时读取这些整数。然而，给定的批次顺序在负载均衡方面可能任意地次优——例如，我们可能有较短的方形预填充后跟长上下文解码。为改善这一点，我们可以对批次强制实施LPT顺序，通过启动一个预处理内核按最大每工作块执行时间对批次进行排序，写出虚拟批次索引到实际批次索引的额外元数据映射，随后注意力内核按排序顺序读回并遍历批次。这些元数据可以被缓存，因此排序不会带来性能损耗。
